\chapter{Assignable Expressions\label{chap:AssignableExpressions}}

\hypertarget{def-assignableexpression}{}
We refer to expressions that may appear on the left-hand side of assignments as \assignableexpressions.
These refer to storage that needs to be updated.

\ExampleDef{Assignable Expressions}
\listingref{AssignableExpressions} shows examples of some assignable expressions.
\ASLListing{Examples of assignable expressions}{AssignableExpressions}{\definitiontests/AssignableExpressions.asl}

Recall that when the \assignableexpression{} is a call to a setter
followed by a slice or field access, it is desugared via \ASTRuleRef{DesugarSetter}
to satisfy \\ \SyntacticSugarRef{SetterFieldAssignment}.

\ChapterOutline
\begin{itemize}
  \item \FormalRelationsRef{Assignable Expressions} introduces the formal relations for
        \assignableexpressions;
  \item \SyntaxRef{Assignable Expressions} defines the syntax for \assignableexpressions;
  \item \AbstractSyntaxRef{Assignable Expressions} defines the functions needed to build
        the AST for \assignableexpressions{}. These make use of desugaring functions,
        defined in \secref{AssignableExpressionsDesugaring};
  \item \secref{DiscardingAssignmentExpressions} defines discarding assignment expressions;
  \item \secref{VariableAssignmentExpressions} defines variable assignment expressions;
  \item \secref{MultiAssignmentExpressions} defines multi-assignment expressions;
  \item \secref{ArrayAssignmentExpressions} defines array assignment expressions;
  \item \secref{BitvectorSliceAssignmentExpressions} defines bitvector slice assignment expressions;
  \item \secref{StructuredTypeFieldAssignmentExpressions} defines structured type field assignment expressions;
  \item \secref{StructuredTypeMultiFieldAssignmentExpressions} defines structured type multi-field assignment expressions;
  \item \secref{BitfieldAssignableExpressions} defines bitfield assignment expressions.
\end{itemize}

%%%%%%%%%%%%%%%%%%%%%%%%%%%%%%%%%%%%%%%%%%%%%%%%%%%%%%%%%%%%%%%%%%%%%%%%%%%%%%%%
\FormalRelationsDef{Assignable Expressions}
\paragraph{Syntax:} Assignable expressions are grammatically derived from $\Nlexpr$.

\paragraph{Abstract Syntax:} Assignable expressions are derived in the abstract syntax from $\lexpr$
  and generated by $\buildlexpr$.

\paragraph{Typing:} Assignable expressions are annotated by $\annotatelexpr$.
\RenderRelation{annotate_lexpr}
\paragraph{Semantics:} Assignable expressions are evaluated by $\evallexpr$.
\RenderRelation{eval_lexpr}
\paragraph{A Convention for Assignable Expressions Semantics Rules:}
In this chapter, variables containing $\vm$ range over $\nativevalue\cartimes\XGraphs$
while variables in which $\vm$ is replaced with $\vv$ correspond to their value component.
For example, $\rmarray\reverseeqdef(\rvarray,\vgtwo)$ and $\mindex\reverseeqdef(\vindex, \vgone)$.

\paragraph{Viewing Assignable Expressions as Right-hand-side Expressions:}
Some of the typing rules and semantics rules require viewing \assignableexpressions\
as \rhsexpressions.
The correspondence is given by the function $\torexpr(\lexpr) \rightarrow \expr$, defined in \secref{LeftToRight}.
%
For example, \SemanticsRuleRef{LESetField}
needs to evaluate the record subexpression $\rerecord$, which is an \assignableexpression.
To achieve this, $\torexpr(\record)$ is used to obtain an \rhsexpression, which then allows
using $\texttt{eval\_expr}$ to evaluate it.

%%%%%%%%%%%%%%%%%%%%%%%%%%%%%%%%%%%%%%%%%%%%%%%%%%%%%%%%%%%%%%%%%%%%%%%%%%%%%%%%
\SyntaxDef{Assignable Expressions}
\begin{flalign*}
\Nlexpr \derives\
   & \Tminus &\\
|\ & \Nbasiclexpr &\\
|\ & \Tlpar \parsesep \Clisttwo{\Ndiscardorbasiclexpr} \parsesep \Trpar &\\
|\ & \Tidentifier \parsesep \Tdot \parsesep \Tlbracket \parsesep \Clisttwo{\Tidentifier} \parsesep \Trbracket &\\
|\ & \Tidentifier \parsesep \Tdot \parsesep \Tlpar \parsesep \Clisttwo{\Ndiscardoridentifier} \parsesep \Trpar &
\end{flalign*}

\begin{flalign*}
\Nbasiclexpr \derives\
   & \Tidentifier \parsesep \Naccess &\\
|\ & \Tidentifier \parsesep \Naccess \parsesep \Nslices &
\end{flalign*}

\begin{flalign*}
\Naccess \derives\
   & \emptysentence  &\\
|\ & \Tdot \parsesep \Tidentifier \parsesep \Naccess &\\
|\ & \Tllbracket \parsesep \Nexpr \parsesep \Trrbracket \parsesep \Naccess &\\
\end{flalign*}

\begin{flalign*}
\Ndiscardorbasiclexpr \derives\ & \Tminus \;|\; \Nbasiclexpr &
\end{flalign*}

\begin{flalign*}
\Ndiscardoridentifier \derives \ & \Tminus \;|\; \Tidentifier &
\end{flalign*}

\SyntacticSugarDef{MultiFieldAssignment}
The syntax
\texttt{x.(f1, ..., fk) = (v1, ..., vk);} is
\syntacticsugar{} for \texttt{x.f1=v1; ... x.fk=vk;},
modulo assignments to \discardvarstr{}.
See \ExampleRef{Desugaring a left-hand side fields tuple}.
It is \desugared{} by \ASTRuleRef{DesugarLHSFieldsTuple}.

%%%%%%%%%%%%%%%%%%%%%%%%%%%%%%%%%%%%%%%%%%%%%%%%%%%%%%%%%%%%%%%%%%%%%%%%%%%%%%%%
\AbstractSyntaxDef{Assignable Expressions}

We first define $\lhsaccess$, which we use in this section as an intermediate representation
between some syntax forms of \assignableexpressions{} and their corresponding abstract syntax.
In particular, rather than directly building the abstract syntax for these \assignableexpressions{},
we first build intermediate structures containing $\lhsaccess$,
which we then desugar into abstract syntax in \secref{AssignableExpressionsDesugaring}.

\hypertarget{ast-lhsaccess}{}
\begin{flalign*}
\lhsaccess \derives\ &
{
\left\{
  \begin{array}{rcl}
    \lhsaccessaccess &:& \fieldorarrayaccess^{*}, \\
    \lhsaccessslices &:& \slice^{*}
\end{array}
\right\}
} &
\hypertarget{ast-fieldorarrayaccess}{}
\hypertarget{ast-fieldaccess}{}
\hypertarget{ast-arrayaccess}{}
\\
\fieldorarrayaccess \derives\ & \FieldAccess(\Identifier)\ |\ \ArrayAccess(\expr)
\end{flalign*}

\ASTRuleDef{LExpr}
\hypertarget{build-lexpr}{}
The function
\[
  \buildlexpr(\overname{\parsenode{\Nlexpr}}{\vparsednode}) \;\aslto\; \overname{\lexpr}{\vastnode}
\]
transforms a parse node $\vparsednode$ into an AST node $\vastnode$.

\begin{mathpar}
\inferrule[discard]{}{
  \buildlexpr(\Nlexpr(\Tminus)) \astarrow \overname{\LEDiscard}{\vastnode}
}
\end{mathpar}

\begin{mathpar}
\inferrule[basic\_lexpr]{
  \desugarlhsaccess(\astof{\vbasiclexpr}) \astarrow \vastnode
}{
  \buildlexpr(\Nlexpr(\punnode{\Nbasiclexpr})) \astarrow \vastnode
}
\end{mathpar}

\begin{mathpar}
\inferrule[multi\_lexpr]{
  \buildclist[\builddiscardorbasiclexpr](\vlexprs) \astarrow \vlexprasts \\
  \desugarlhstuple(\vlexprasts) \astarrow \vastnode
}{
  {
  \begin{array}{r}
  \buildlexpr(\Nlexpr(\Tlpar, \namednode{\vlexprs}{\Clisttwo{\Ndiscardorbasiclexpr}}, \Trpar)) \\
  \astarrow \vastnode
  \end{array}
  }
}
\end{mathpar}

\begin{mathpar}
\inferrule[concat\_fields]{
  \buildclist[\buildidentity](\vfields) \astarrow \vfieldasts
}{
  {
  \begin{array}{r}
  \buildlexpr(\Nlexpr(\Tidentifier(\id), \Tdot, \Tlbracket, \namednode{\vfields}{\Clisttwo{\Tidentifier}}, \Trbracket)) \astarrow \\
  \overname{\LESetFields (\LEVar(\id), \vfieldasts)}{\vastnode}
  \end{array}
  }
}
\end{mathpar}

\begin{mathpar}
\inferrule[tuple\_fields]{
  \buildclist[\builddiscardoridentifier](\vids) \astarrow \astversion{\vids} \\
  \desugarlhsfieldstuple(\id, \astversion{\vids}) \astarrow \vastnode
}{
  {
  \begin{array}{r}
  \buildlexpr(\Nlexpr(\Tidentifier(\id), \Tdot, \Tlpar, \namednode{\vids}{\Clisttwo{\Ndiscardoridentifier}}, \Trpar))\\ \astarrow
  \vastnode
  \end{array}
  }
}
\end{mathpar}

\ASTRuleDef{BasicLexpr}
\hypertarget{build-basiclexpr}{}
The function
\[
  \buildbasiclexpr(\overname{\parsenode{\Nbasiclexpr}}{\vparsednode}) \;\aslto\; (\overname{\Identifier}{\vbase} \aslsep \overname{\lhsaccess}{\vlhsaccess})
\]
transforms a parse node $\vparsednode$ into a pair of AST nodes, $\vbase$ and $\vlhsaccess$.

\begin{mathpar}
\inferrule[no\_slices]{
  {
  \vlhsaccess \eqdef
    \left\{
      \begin{array}{rcl}
        \lhsaccessaccess &:& \astof{\vaccess},\\
        \lhsaccessslices &:& \emptylist
      \end{array}
    \right\}
  }
}{
  \buildbasiclexpr(\Nbasiclexpr(\Tidentifier(\vbase),\punnode{\Naccess})) \astarrow
  (\vbase, \vlhsaccess)
}
\end{mathpar}

\begin{mathpar}
\inferrule[slices]{
  {
  \vlhsaccess \eqdef
    \left\{
      \begin{array}{rcl}
              \lhsaccessaccess &:& \astof{\vaccess},\\
              \lhsaccessslices &:& \astof{\vslices}
      \end{array}
    \right\}
  }
}{
  \buildbasiclexpr(\Nbasiclexpr(\Tidentifier(\vbase),\punnode{\Naccess}, \punnode{\Nslices})) \astarrow
  (\vbase, \vlhsaccess)
}
\end{mathpar}

\ASTRuleDef{Access}

\hypertarget{build-access}{}
The function
\[
  \buildaccess(\overname{\parsenode{\Naccess}}{\vparsednode}) \;\aslto
  \overname{\KleeneStar{\fieldorarrayaccess}}{\vastnode}
\]
transforms a parse node $\vparsednode$ into an AST node $\vastnode$.

\begin{mathpar}
\inferrule[empty]{}{
  \buildaccess(\Naccess(\emptysentence)) \astarrow \overname{\emptylist}{\vastnode}
}
\end{mathpar}

\begin{mathpar}
\inferrule[field\_access]{}{
  \buildaccess(\Naccess(\Tdot, \Tidentifier(\vfield), \punnode{\Naccess})) \astarrow
  \overname{[\FieldAccess(\vfield)] \concat \astof{\vaccess}}{\vastnode}
}
\end{mathpar}

\begin{mathpar}
\inferrule[array\_access]{}{
  \buildaccess(\Naccess(\Tllbracket, \punnode{\Nexpr}, \Trrbracket, \punnode{\Naccess})) \astarrow
  \overname{[\ArrayAccess(\astof{\vexpr})] \concat \astof{\vaccess}}{\vastnode}
}
\end{mathpar}

\ASTRuleDef{DiscardOrBasicLexpr}
\hypertarget{build-discardorbasiclexpr}{}
The function
\[
\begin{array}{r}
  \builddiscardorbasiclexpr(\overname{\parsenode{\Ndiscardorbasiclexpr}}{\vparsednode}) \;\aslto\\
  \overname{\Option{\lhsaccess}}{\vastnode}
\end{array}
\]
transforms a parse node $\vparsednode$ into an AST node $\vastnode$.

\begin{mathpar}
\inferrule[discard]{}{
  \builddiscardorbasiclexpr(\Ndiscardorbasiclexpr(\Tminus)) \astarrow \overname{\None}{\vastnode}
}
\end{mathpar}

\begin{mathpar}
\inferrule[basic]{}{
  {
    \begin{array}{r}
      \builddiscardorbasiclexpr(\Ndiscardorbasiclexpr(\punnode{\Nbasiclexpr})) \astarrow \\
      \overname{\some{ \astof{\vbasiclexpr} }}{\vastnode}
    \end{array}
  }
}
\end{mathpar}

\ASTRuleDef{DiscardOrIdentifier}
\hypertarget{build-discardoridentifier}{}
The function
\[
\begin{array}{r}
  \builddiscardoridentifier(\overname{\parsenode{\Ndiscardoridentifier}}{\vparsednode}) \;\aslto
  \overname{\Option{\Identifier}}{\vastnode}
\end{array}
\]
transforms a parse node $\vparsednode$ into an AST node $\vastnode$.

\begin{mathpar}
\inferrule[none]{}{
  \builddiscardoridentifier(\Ndiscardoridentifier(\Tminus)) \astarrow \overname{\None}{\vastnode}
}
\end{mathpar}

\begin{mathpar}
\inferrule[some]{}{
  {
    \begin{array}{r}
      \builddiscardoridentifier(\Ndiscardoridentifier(\Tidentifier(\id))) \astarrow
      \overname{\some{ \id }}{\vastnode}
    \end{array}
  }
}
\end{mathpar}

\subsection{Desugaring Assignable Expressions\label{sec:AssignableExpressionsDesugaring}}

This section defines three desugaring relations which produce assignable expression abstract syntax, that is, $\lexpr$.
They are used to build $\lexpr$ abstract syntax.
\begin{itemize}
  \item $\desugarlhsaccess$, which desugars a pair consisting of an $\Identifier$ and an $\lhsaccess$ into an $\lexpr$.
  \item $\desugarlhstuple$, which desugars a tuple of optional $\lhsaccess$ elements into an $\lexpr$.
        This represents a multi-assignment of a tuple value, where $\None$ means that element of the tuple is discarded.
  \item $\desugarlhsfieldstuple$, which desugars a multi-assignment of a tuple value to multiple fields of an identifier.
\end{itemize}

\ASTRuleDef{DesugarLHSAccess}
\hypertarget{def-desugarlhsaccess}{}
The function
\[
  \desugarlhsaccess(\overname{\Identifier}{\name}, \overname{\lhsaccess}{\vlhsaccess}) \;\aslto\; \overname{\lexpr}{\vlexpr}
\]
transforms an $\lhsaccess$ on an identifier $\name$ into an AST node $\vlexpr$.

\ExampleDef{Desugaring a left-hand side access}
\listingref{DesugarLHSAccess} shows an example of desugaring an $\lhsaccess$.
\ASLListing{Desugaring a left-hand side access}{DesugarLHSAccess}{\syntaxtests/ASTRule.DesugarLHSAccess.asl}

\begin{mathpar}
\inferrule{
  \vlexprs_0 \eqdef \LEVar(\name) \\
  {
  i \in 1..|\vaccess|: \vlexprs_i \eqdef
    \begin{cases}
      \LESetField(\vlexprs_{i-1}, \vfield) \\
        \qquad \text{if } \vaccess_i = \FieldAccess(\vfield) \\
      \LESetArray(\vlexprs_{i-1}, \vexpr) \\
        \qquad \text{if } \vaccess_i = \ArrayAccess(\vexpr) \\
    \end{cases}
  } \\
  \vlexpr \eqdef
  \ifthenelseop[H]{\vslices = \emptylist}{\vlexprs_{|\vaccess|}}{\LESlice(\vlexprs_{|\vaccess|}, \vslices)}
}{
  \desugarlhsaccess\overname{
    \left\{
      {
        \begin{array}{rcl}
          \lhsaccessaccess &:& \vaccess, \\
          \lhsaccessslices &:& \vslices
        \end{array}
      }
    \right\}
  }{\vlhsaccess} \astarrow \vlexpr
}
\end{mathpar}

\ASTRuleDef{DesugarLHSTuple}
\hypertarget{def-desugarlhstuple}{}
The function
\[
  \desugarlhstuple(\overname{\KleeneStar{\Option{\lhsaccess}}}{\vlhsaccessopts}) \;\aslto\; \overname{\lexpr}{\vlexpr}
\]
transforms a list of \optionalterm{} $\lhsaccess$ elements into an AST node $\vlexpr$.

\ExampleDef{Desugaring a left-hand side tuple}
\listingref{DesugarLHSTuple} shows an example of desugaring a left-hand side tuple.
\ASLListing{Desugaring a left-hand side tuple}{DesugarLHSTuple}{\syntaxtests/ASTRule.DesugarLHSTuple.asl}

\begin{mathpar}
\inferrule{
  i \in 1..|\vlhsaccessopts|: \desugarlhsaccessopt(\vlhsaccessopt_i) \astarrow \vlexpr_i \\
  \vlexprs \eqdef [i \in 1..|\vlhsaccessopts|: \vlexpr_i]
}{
  \desugarlhstuple(\vlhsaccessopts) \astarrow \overname{\LEDestructuring(\vlexprs)}{\vlexpr}
}
\end{mathpar}

\ASTRuleDef{DesugarLHSAccessOpt}
\hypertarget{def-desugarlhsaccessopt}{}
The helper AST function
\[
    \desugarlhsaccessopt(\overname{\Option{\lhsaccess}}{\vlhsaccessopt}) \;\aslto\; \overname{\lexpr}{\vlexpr}
\]
is defined via the following rules:

\begin{mathpar}
\inferrule[none]{}
{
  \desugarlhsaccessopt(\overname{\None}{\vlhsaccessopt}) \astarrow \overname{\LEDiscard}{\vlexpr}
}
\end{mathpar}

\begin{mathpar}
\inferrule[some]{
  \desugarlhsaccess(\vlhsaccess) \astarrow \vlexpr
}{
  \desugarlhsaccessopt(\overname{\some{\vlhsaccess}}{\vlhsaccessopt}) \astarrow \vlexpr
}
\end{mathpar}

\ASTRuleDef{DesugarLHSFieldsTuple}
\hypertarget{def-desugarlhsfieldstuple}{}
The function
\[
  \desugarlhsfieldstuple(\overname{\Identifier}{\id} \aslsep \overname{\KleeneStar{\Option{\Identifier}}}{\fieldopts}) \;\aslto\; \overname{\lexpr}{\vlexpr}
\]
transforms an assignment to a tuple of fields $\fields$ of variable $\id$ into an AST node $\vlexpr$.

\ExampleDef{Desugaring a left-hand side fields tuple}
\listingref{DesugarLHSFieldsTuple} shows an example of desugaring a left-hand side fields tuple.
\ASLListing{Desugaring a left-hand side fields tuple}{DesugarLHSFieldsTuple}{\syntaxtests/ASTRule.DesugarLHSFieldsTuple.asl}

\begin{mathpar}
\inferrule{
  \fields \eqdef \filteroptionlist{\fieldopts}\\
  \checknoduplicates(\fields) \typearrow \True \OrTypeError \hva\\
  i \in 1..|\fieldopts|: \desugarlhsfieldopt(\id, \fieldopts[i]) \astarrow \vlexpr_i \\
  \vlexprs \eqdef [i \in 1..|\fieldopts|: \vlexpr_i]
}{
  \desugarlhsfieldstuple(\id, \fieldopts) \astarrow \overname{\LEDestructuring(\vlexprs)}{\vlexpr}
}
\end{mathpar}

\ASTRuleDef{DesugarLHSFieldOpt}
\hypertarget{def-desugarlhsfieldopt}{}
The helper AST function
\[
    \desugarlhsfieldopt(\overname{\Identifier}{\id} \aslsep \overname{\Option{\Identifier}}{\fieldopt}) \;\aslto\; \overname{\lexpr}{\vlexpr}
\]
is defined via the following rules:
\begin{mathpar}
\inferrule[none]{}
{
  \desugarlhsfieldopt(\id, \overname{\None}{\fieldopt}) \astarrow \overname{\LEDiscard}{\vlexpr}
}
\end{mathpar}

\begin{mathpar}
\inferrule[some]{}{
  \desugarlhsfieldopt(\id, \overname{\some{\vfield}}{\fieldopt}) \astarrow \overname{\LESetField(\LEVar(\id), \vfield)}{\vlexpr}
}
\end{mathpar}

\hypertarget{def-discardlexprterm}{}
\section{Discarding Assignment Expressions\label{sec:DiscardingAssignmentExpressions}}
\subsection{Abstract Syntax}
\RenderTypes[remove_hypertargets]{lexpr_discard}


\subsection{Typing}
\TypingRuleDef{LEDiscard}
\ExampleDef{Well-typed Discarding Assignments}
All discarding assignments in \listingref{typing-lediscard} are well-typed.
\ASLListing{Well-typed discarding assignments}{typing-lediscard}{\typingtests/TypingRule.LEDiscard.asl}


\RenderProseAndFormally{annotate_lexpr_LEDiscard}
\CodeSubsection{\LEDiscardBegin}{\LEDiscardEnd}{../Typing.ml}

\subsection{Semantics}
\SemanticsRuleDef{LEDiscard}
\ExampleDef{Discarding Assignments}
In \listingref{semantics-lediscard}, the assignment \texttt{- = 42;} does not affect the environment.
\ASLListing{Assignment to \texttt{-}}{semantics-lediscard}{\semanticstests/SemanticsRule.LEDiscard.asl}


\RenderProseAndFormally{eval_lexpr_LEDiscard}
\CodeSubsection{\EvalLEDiscardBegin}{\EvalLEDiscardEnd}{../Interpreter.ml}

\hypertarget{def-varlexprterm}{}
\section{Variable Assignment Expressions\label{sec:VariableAssignmentExpressions}}
\subsection{Abstract Syntax}
\RenderTypes[remove_hypertargets]{lexpr_var}


\subsection{Typing}
\TypingRuleDef{LEVar}
\ExampleDef{Variable Assignments}
In \listingref{typing-levar}, all variable assignments are well-typed.
\ASLListing{Well-typed variable assignments}{typing-levar}{\typingtests/TypingRule.LEVar.asl}

In \listingref{typing-levar-undefined}, the assignment to \verb|x| is ill-typed,
since \verb|x| is not defined as either a local storage element or as a global storage element.
\ASLListing{An ill-typed variable assignment}{typing-levar-undefined}{\typingtests/TypingRule.LEVar.undefined.asl}


\RenderProseAndFormally{annotate_lexpr_LEVar}


\CodeSubsection{\LEVarBegin}{\LEVarEnd}{../Typing.ml}

\subsection{Semantics}
\SemanticsRuleDef{LEVar}
\ExampleDef{Local Variable}
In \listingref{semantics-levar}, \SemanticsRuleRef{LEVar} is (only) used
to assign the value $42$ to the left-hand-side expression
\texttt{x} within \texttt{x = 42;}.

\ASLListing{Semantics of a left-hand-side variable expression}{semantics-levar}{\semanticstests/SemanticsRule.LELocalVar.asl}

\ExampleDef{Global Variable}
In \listingref{semantics-leglobalvar},
\SemanticsRuleRef{LEVar} is (only) used to assign the value $42$ to the left-hand-side expression
\texttt{x} within \texttt{x = 42;}.
\ASLListing{Assignment to a global variable}{semantics-leglobalvar}{\semanticstests/SemanticsRule.LEGlobalVar.asl}


\RenderProseAndFormally{eval_lexpr_LEVar}
\CodeSubsection{\EvalLEVarBegin}{\EvalLEVarEnd}{../Interpreter.ml}

\section{Multi-assignment Expressions\label{sec:MultiAssignmentExpressions}}
\hypertarget{def-destructuringlexprterm}{}
\ASLListing{Assignment to multiple left-hand-side expressions}{semantics-ledestructuring}{\semanticstests/SemanticsRule.LEDestructuring.asl}

\subsection{Abstract Syntax}
\RenderTypes[remove_hypertargets]{lexpr_destructuring}


\subsection{Typing}
\TypingRuleDef{LEDestructuring}
\ExampleDef{Well-typed Multi-variable Assignment}
In \listingref{semantics-ledestructuring}, the multi-assignment
\verb|(x, y) = (3, 42)| is well-typed.


\RenderProseAndFormally{annotate_lexpr_LEDestructuring}
\CodeSubsection{\LEDestructuringBegin}{\LEDestructuringEnd}{../Typing.ml}

\subsection{Semantics}
\SemanticsRuleDef{LEDestructuring}
\ExampleDef{Multi-assignments}
In \listingref{semantics-ledestructuring}, the multi-assignment \verb|(x, y) = (3, 42)| binds
\texttt{x} to $\nvint(3)$ and \texttt{y} to
$\nvint(42)$ in the environment where \texttt{x} is bound to $\nvint(42)$ and
\texttt{y} is bound to $\nvint(3)$.


\RenderProseAndFormally{eval_lexpr_LEDestructuring}
\CodeSubsection{\EvalLEDestructuringBegin}{\EvalLEDestructuringEnd}{../Interpreter.ml}

\SemanticsRuleDef{LEMultiAssign}
\RenderRelation{multi_assign}
See \ExampleRef{Multi-assignments}.

\RenderRule{multi_assign}

Notice that this rule is only defined when the lists $\vlelist$ and $\vmlist$ have the same length.
To see this, notice that
to form a derivation tree, we must employ the \textsc{non\_empty} case, which ensures both lists
have at least one element and shortens the lengths of both lists by one,
until both lists become empty,
which is when the \textsc{empty} axiom case is used.

\SemanticsRuleDef{ProtectedLEMultiAssign}
\RenderRelation{protected_multi_assign}
See \ExampleRef{Evaluation of Multi-variable Assignment from Subprogram Calls}
and \ExampleRef{Multi-assignments}.

\ProseParagraph
\OneApplies
\begin{itemize}
  \item \AllApplyCase{same\_length}
  \begin{itemize}
    \item $\vlelist$ and $\vmlist$ have the same length;
    \item applying $\multiassign$ to $\env$, $\vlelist$, and $\vmlist$ yields\\
          $\ResultLexpr(\vg, \newenv)$\ProseOrAbnormal;
    \item \Proseeqdef{$\newg$}{$\vg$}.
  \end{itemize}

  \item \AllApplyCase{different\_length}
  \begin{itemize}
    \item $\vlelist$ and $\vmlist$ have different lengths;
    \item $\vmlist$ is a singleton list containing the value $(\vtuplevalue, \vg)$;
    \item applying $\getindex$ to $\vi$, and $\vtuplevalue$ for each index
          $\vi$ in the list of indices for $\vlelist$ yields $\vtupleelements[\vi]$\ProseOrError;
    \item \Proseeqdef{$\nmlist$}{the pair $\vtupleelements[\vi]$ and $\vg$ for each index $\vi$ in the list of indices for $\vlelist$};
    \item applying $\multiassign$ to $\env$, $\vlelist$, and $\nmlist$ yields\\
          $\ResultLexpr(\newg, \newenv)$\ProseOrAbnormal.
  \end{itemize}
\end{itemize}

\FormallyParagraph
\RenderRule{protected_multi_assign}

\hypertarget{def-setarraylexprterm}{}
\section{Array Assignment Expressions\label{sec:ArrayAssignmentExpressions}}
This section details the syntax, abstract syntax, semantics, and typing of array write expressions.
In the untyped AST, a write to either an integer-indexed array or an enumeration-indexed array is represented
the same way. The type system infers the kind of array and outputs a typed AST node differentiating
the two kinds of arrays, either a $\LESetArray$ or a $\LESetEnumArray$, via \TypingRuleRef{LESetArray}.
The semantics utilizes a rule matching the corresponding type of array ---
\SemanticsRuleRef{LESetArray} for integer-indexed arrays and
\SemanticsRuleRef{LESetEnumArray} for enumeration-indexed arrays.

\subsection{Abstract Syntax}
\RenderTypes[remove_hypertargets]{lexpr_setarray_and_typed}


\subsection{Typing}
\TypingRuleDef{LESetArray}
The assignable expressions in
\listingref{semantics-lesetarray} and \listingref{semantics-lesetenumarray}
are well-typed.


\RenderProseAndFormally{annotate_lexpr_LESetArray}
\CodeSubsection{\LESetArrayBegin}{\LESetArrayEnd}{../Typing.ml}

\TypingRuleDef{AnnotateSetArray}
\RenderRelation{annotate_set_array}
See \listingref{semantics-lesetarray} and \listingref{semantics-lesetenumarray}
for examples of well-typed assignable array expressions.


\RenderProseAndFormally{annotate_set_array}

\subsection{Semantics}
\SemanticsRuleDef{LESetArray}
\ExampleDef{Integer-indexed Array Update Assignments}
In \listingref{semantics-lesetarray}, the assignment \verb|my_array[[3]] = 53;| binds the third element
of \texttt{my\_array} to the value \texttt{53}.
\ASLListing{Assignment to an integer-indexed array cell}{semantics-lesetarray}{\semanticstests/SemanticsRule.LESetArray.asl}


\RenderProseAndFormally{eval_lexpr_LESetArray}
\CodeSubsection{\EvalLESetArrayBegin}{\EvalLESetArrayEnd}{../Interpreter.ml}

\subsubsection{Comments}
We note that the index is guaranteed by the typechecker to be within the array bounds
via \TypingRuleRef{LESetArray}.

See \SyntacticSugarRef{SetterFieldAssignment}.

\SemanticsRuleDef{LESetEnumArray}
\ExampleDef{Enumeration-indexed Array Update Assignments}
In \listingref{semantics-lesetenumarray}, the assignment \verb|my_array[[RED]] = 53;| binds the \verb|RED|
cell of
\texttt{my\_array} to the value \texttt{53}.
\ASLListing{Assignment to an enumeration-indexed array cell}{semantics-lesetenumarray}{\semanticstests/SemanticsRule.LESetEnumArray.asl}


\RenderProseAndFormally{eval_lexpr_LESetEnumArray}
\CodeSubsection{\EvalLESetEnumArrayBegin}{\EvalLESetEnumArrayEnd}{../Interpreter.ml}

\hypertarget{def-slicelexprterm}{}
\section{Bitvector Slice Assignment Expressions\label{sec:BitvectorSliceAssignmentExpressions}}
\ASLListing{Assignable slice expressions}{semantics-leslice}{\semanticstests/SemanticsRule.LESlice.asl}

\subsection{Abstract Syntax}
\RenderTypes[remove_hypertargets]{lexpr_slice}


\subsection{Typing}
\TypingRuleDef{LESlice}
\ExampleDef{Typing Assignable Slice Expression}
In \listingref{semantics-leslice}, the assignable slice expression \verb|x[3:0, 7:6]| is well-typed.

In \listingref{typing-leslice-bad}, the assignable slice expression \verb|x[3:0, 3]| is ill-typed,
since the slice \verb|3:0| overlaps with the slice \verb|3| at position $3$.
The specification is ill-typed, even though both slices assign \verb|0| to the bit at position $3$.
\ASLListing{An ill-typed assignable slice expression}{typing-leslice-bad}{\typingtests/TypingRule.LESlice.bad.asl}


\RenderProseAndFormally{annotate_lexpr_LESlice}
\CodeSubsection{\LESliceBegin}{\LESliceEnd}{../Typing.ml}

\TypingRuleDef{CheckDisjointSlices}
\RenderRelation{check_disjoint_slices}
See \ExampleRef{Typing Assignable Slice Expression}.


\RenderProseAndFormally{check_disjoint_slices}
\CodeSubsection{\CheckDisjointSlicesBegin}{\CheckDisjointSlicesEnd}{../Typing.ml}

\subsection{Semantics}
\SemanticsRuleDef{LESlice}
\ExampleDef{Slice Assignments}
In \listingref{semantics-leslice}, the assignment\\ \texttt{x[3:0, 7:6] = '000000';} binds
\texttt{x} to $\nvbitvector(00110000)$
via the rule \SemanticsRuleRef{LESlice}
in the environment where \texttt{x} is bound to $\nvbitvector(11111111)$.


\RenderProseAndFormally{eval_lexpr_LESlice}
\CodeSubsection{\EvalLESliceBegin}{\EvalLESliceEnd}{../Interpreter.ml}
\subsubsection{Comments}
See \SyntacticSugarRef{SetterFieldAssignment}.

\SemanticsRuleDef{CheckNonOverlappingSlices}
\RenderRelation{check_non_overlapping_slices}
\ExampleDef{Checking Slices for Overlaps}
In \listingref{semantics-checknonoverlappingslices},
the slices \verb|N-1:N-2| and \verb|0| do not overlap, and\\ $\checknonoverlappingslices$ yields $\True$.
\ASLListing{Non-overlapping slices}{semantics-checknonoverlappingslices}{\semanticstests/SemanticsRule.CheckNonOverlappingSlices.asl}

In \listingref{semantics-checknonoverlappingslices-bad},
the slices \verb|N-1:N-2| and \verb|0| overlap, and $\checknonoverlappingslices$
yields a \dynamicerrorterm, even though the assignment \verb|y[N-1:N-2, 0] = '111';|
writes \verb|1| to the common position $0$.
\ASLListing{Overlapping slices}{semantics-checknonoverlappingslices-bad}{\semanticstests/SemanticsRule.CheckNonOverlappingSlices.bad.asl}


\RenderProseAndFormally{check_non_overlapping_slices}

\SemanticsRuleDef{CheckTwoRangesNonOverlapping}
\RenderRelation{check_two_ranges_non_overlapping}
See \ExampleRef{Checking Slices for Overlaps}.


\RenderProseAndFormally{check_two_ranges_non_overlapping}

\hypertarget{def-setfieldlexprterm}{}
\section{Structured Type Field Assignment Expressions\label{sec:StructuredTypeFieldAssignmentExpressions}}
\subsection{Abstract Syntax}
\RenderTypes[remove_hypertargets]{lexpr_setfield}

\subsection{Typing}

\TypingRuleDef{LESetStructuredField}
\ExampleDef{Typing a Structured Type Field Assignment}
In \listingref{typing-lestructuredfield}, all field assignment statements
are well-typed.
\ASLListing{Typing a structured type field assignment}{typing-lestructuredfield}{\typingtests/TypingRule.LESetStructuredField.asl}


\RenderProseAndFormally{annotate_lexpr_LESetStructuredField}
\CodeSubsection{\LESetStructuredFieldBegin}{\LESetStructuredFieldEnd}{../Typing.ml}

\TypingRuleDef{LESetCollectionField}

\ExampleDef{Typing Collection Field Assignable Expressions}
All of the collection field assignable expressions in
\listingref{typing-lesetcollectionfields} are well-typed.


\RenderProseAndFormally{annotate_lexpr_LESetCollectionField}
\CodeSubsection{\LESetStructuredFieldBegin}{\LESetStructuredFieldEnd}{../Typing.ml}

\TypingRuleDef{LESetBadField}
\ExampleDef{Assigning a Field in an Inappropriate Type}
In \listingref{typing-lesetbadfield}, the statement \verb|x.RED = 42;| is ill-typed,
since \verb|x| is not a \bitvectortypeterm{} nor a \structuredtypeterm.
\ASLListing{Assigning a field in an inappropriate type}{typing-lesetbadfield}{\typingtests/TypingRule.LESetBadField.asl}

See \RequirementRef{TupleImmutability} for an example of assigning to a tuple type.


\RenderProseAndFormally{annotate_lexpr_LESetBadField_error}
\CodeSubsection{\LESetBadFieldBegin}{\LESetBadFieldEnd}{../Typing.ml}

\SemanticsRuleDef{LESetField}

\ExampleDef{Field Assignment}
In \listingref{semantics-lesetfield}, the assignment
\verb|my_record.a = 42;| binds \\ \verb|my_record| to \verb|{a: 42, b: 100}|
in the environment where \verb|my_record| is bound to \\
\verb|{a: 3, b: 100}|.

\ASLListing{Assignment to a field}{semantics-lesetfield}{\semanticstests/SemanticsRule.LESetField.asl}


\RenderProseAndFormally{eval_lexpr_LESetField}
\CodeSubsection{\EvalLESetFieldBegin}{\EvalLESetFieldEnd}{../Interpreter.ml}

\subsubsection{Comments}
We note that the typechecker guarantees that $\fieldname$ exists in the record given by $\record$
via \TypingRuleRef{LESetStructuredField}.

See \SyntacticSugarRef{SetterFieldAssignment}.

\section{Structured Type Multi-field Assignment Expressions\label{sec:StructuredTypeMultiFieldAssignmentExpressions}}

\ASLListing{Multi-field assignment expression}{lesetfields}{\typingtests/TypingRule.LESetFields.asl}

\subsection{Abstract Syntax}
\RenderTypes[remove_hypertargets]{lexpr_setfields}

\subsection{Typing}
\TypingRuleDef{LESetFields}
All multi-field assignment expressions in \listingref{lesetfields} are well-typed.


\RenderProseAndFormally{annotate_lexpr_LESetFields}
\CodeSubsection{\LESetFieldsBegin}{\LESetFieldsEnd}{../Typing.ml}

\TypingRuleDef{FoldBitvectorFields}
\RenderRelation{fold_bitvector_fields}
\ExampleDef{Obtaining the Total Width and Ranges of Bitvector-typed Fields}
In \listingref{lesetfields}, applying $\foldbitvectorfields$ to the list of fields
[\verb|data, time, status|] in the statement \verb|x.[status, time, data] = '1' :: Zeros{16} :: Ones{8};|,
yields the total width $25$ and ranges
$(\overtext{17}{\start}, \overtext{8}{\vwidth})$ for \verb|data|,
$(\overtext{1}{\start}, \overtext{16}{\vwidth})$ for \verb|time|,
and $(\overtext{0}{\start}, \overtext{1}{\vwidth})$ for \verb|status|.


\RenderProseAndFormally{fold_bitvector_fields}

\subsection{Semantics}
\SemanticsRuleDef{LESetFields}
The multi-field assignments in \listingref{lesetfields} evaluate without yielding a \dynamicerrorterm.


\RenderProseAndFormally{eval_lexpr_LESetFields}

\SemanticsRuleDef{AssignBitvectorFields}
\RenderRelation{assign_bitvector_fields}
\ExampleDef{Assignment Bitvector-typed Fields}
The statement \verb|y.[data, time, status] = Ones{8} :: Zeros{16} :: '1';|
in \listingref{lesetfields}
assigns the fields \verb|data|, \verb|time|, and \verb|status|,
yielding \\
$\nvbitvector('11111111 0000000000000000 1')$.


\RenderProseAndFormally{assign_bitvector_fields}
\SemanticsRuleDef{LESetCollectionFields}

\ExampleDef{Typing Collection Fields Assignable Expressions}
All of the collection field assignable expressions in
\listingref{typing-lesetcollectionfields} are well-typed.
\ASLListing{Typing collection fields assignment expressions}{typing-lesetcollectionfields}{\typingtests/TypingRule.LESetCollectionFields.asl}


\RenderProseAndFormally{eval_lexpr_LESetCollectionFields}

\section{Bitfield Assignable Expressions\label{sec:BitfieldAssignableExpressions}}
\subsection{Abstract Syntax}
\RenderTypes[remove_hypertargets]{lexpr_slice}


\subsection{Typing}
\TypingRuleDef{LESetBitField}

\ExampleDef{Typing Bitfield Assignable Expressions}
All of the bitfield assignable expressions in \listingref{typing-lesetfield}
are well-typed.
\ASLListing{Typing bitfield assignment expressions}{typing-lesetfield}{\typingtests/TypingRule.LESetField.asl}


\RenderProseAndFormally{annotate_lexpr_LESetField_BitField}
\subsection{Semantics}
The semantics for assigning to individual bitvector bitfields is covered by \SemanticsRuleRef{LESlice}
as the type system transforms the \untypedast\ for assigning to an individual bitfield into an $\LESlice$ \typedast.
